\chapter{Chapter05 — Inner Product Mod Two}
%%% generated by the blueprint pipeline from FABL.Chapter05.InnerProductModTwo
\section{Overview}

Declarations of the Lean module \texttt{FABL.Chapter05.InnerProductModTwo}.

\section{Declarations}

\begin{definition}[Boolean cube block decomposition]
\lean{FABL.f₂CubeBlockEquiv}
\leanok
For each natural number $n$, the equivalence between the Boolean cube of dimension $2n$
and a pair of Boolean cubes of dimension $n$ sends a vector to its first and second
coordinate blocks; its inverse concatenates the two blocks.
\end{definition}

\begin{definition}[Concatenation of Boolean cube blocks]
\lean{FABL.joinF₂CubeBlocks}
\leanok
\uses{FABL.f₂CubeBlockEquiv}
For each natural number $n$ and Boolean cube vectors $x,y$ of dimension $n$, their
joined vector is the vector of dimension $2n$ whose first block is $x$ and whose second
block is $y$.
\end{definition}

\begin{theorem}[Reconstruction of two binary blocks]
\lean{FABL.f₂CubeBlockEquiv_joinF₂CubeBlocks}
\leanok
\uses{FABL.f₂CubeBlockEquiv, FABL.joinF₂CubeBlocks}
\difficulty{1}
For every natural number $n$ and every pair of vectors $x,y\in(\mathbb{F}_2)^n$,
splitting the concatenation of $x$ and $y$ into its two length-$n$ blocks recovers
$(x,y)$.
\end{theorem}

\begin{theorem}[First block of a joined cube]
\lean{FABL.joinF₂CubeBlocks_castAdd}
\leanok
\uses{FABL.joinF₂CubeBlocks}
\difficulty{3}
Let $n$ be a natural number, and let $x,y\in\mathbb{F}_2^n$. Their concatenation in
$\mathbb{F}_2^{2n}$ has value $x_i$ at each coordinate $i$ in the first block.
\end{theorem}

\begin{theorem}[Second block of a joined cube]
\lean{FABL.joinF₂CubeBlocks_natAdd}
\leanok
\uses{FABL.joinF₂CubeBlocks}
\difficulty{3}
Let $n$ be a natural number, and let $x,y\in\mathbb{F}_2^n$. Their concatenation in
$\mathbb{F}_2^{2n}$ has value $y_i$ at coordinate $i$ in the second block.
\end{theorem}

\begin{theorem}[Second block of a concatenated cube vector]
\lean{FABL.joinF₂CubeBlocks_addNat}
\leanok
\uses{FABL.joinF₂CubeBlocks, FABL.joinF₂CubeBlocks_natAdd}
\difficulty{3}
Let $n$ be a natural number, let $x,y\in\mathbb F_2^n$, and let $i$ be an index from $0$
to $n-1$. The vector formed by concatenating $x$ and $y$, evaluated at the position
corresponding to index $i$ in the second block, has value $y_i$.
\end{theorem}

\begin{definition}[Binary inner product bit]
\lean{FABL.innerProductModTwoBit}
\statementsource{Vad12}{
  pdf-bf2ce5422088-p142-b003,
  pdf-bf2ce5422088-p142-b004,
  pdf-bf2ce5422088-p305-b002
}
\statementsource{HILL99}{pdf-130039751e9f-p006-b005}
\statementsource{Analysis-of-Boolean-Functions-by-Ryan-ODonnell}{
  pdf-ffbe9c7b3306-p034-b002,
  pdf-ffbe9c7b3306-p149-b003
}
\label{FABL.innerProductModTwoBit}
\leanok
\uses{FABL.f₂CubeBlockEquiv, FABL.f₂DotProduct, FABL.𝔽₂}
For a vector $z\in\mathbb{F}_2^{2n}$, split $z$ into its first and last $n$ coordinates,
say $u,v\in\mathbb{F}_2^n$. The binary inner product bit of $z$ is $\sum_{i=1}^n u_i
v_i$ computed in $\mathbb{F}_2$.
\end{definition}

\begin{definition}[Sign encoding of the binary inner product]
\lean{FABL.innerProductModTwo}
\statementsource{Vad12}{
  pdf-bf2ce5422088-p142-b003,
  pdf-bf2ce5422088-p142-b004,
  pdf-bf2ce5422088-p305-b002
}
\statementsource{HILL99}{pdf-130039751e9f-p006-b005}
\statementsource{Analysis-of-Boolean-Functions-by-Ryan-ODonnell}{
  pdf-ffbe9c7b3306-p034-b002,
  pdf-ffbe9c7b3306-p149-b003,
  pdf-ffbe9c7b3306-p151-b006
}
\label{FABL.innerProductModTwo}
\leanok
\uses{FABL.innerProductModTwoBit, FABL.realSignEncodedFunction}
For each natural number $n$, the real-valued function on binary vectors of length $2n$
is the real sign encoding of the inner product modulo two of the corresponding bit
vector.
\end{definition}

\begin{definition}[Boolean encoding of the binary inner product]
\lean{FABL.innerProductModTwoBoolean}
\statementsource{Vad12}{
  pdf-bf2ce5422088-p142-b003,
  pdf-bf2ce5422088-p142-b004,
  pdf-bf2ce5422088-p305-b002
}
\statementsource{HILL99}{pdf-130039751e9f-p006-b005}
\statementsource{Analysis-of-Boolean-Functions-by-Ryan-ODonnell}{
  pdf-ffbe9c7b3306-p034-b002,
  pdf-ffbe9c7b3306-p149-b003
}
\label{FABL.innerProductModTwoBoolean}
\leanok
\uses{FABL.BooleanFunction, FABL.binaryCubeSignEquiv, FABL.innerProductModTwoBit, FABL.signEncode}
For each natural number $n$, the Boolean function on sign vectors of length $2n$
converts each sign vector to its corresponding binary vector, computes the inner product
modulo two, and encodes the result as a sign.
\end{definition}

\begin{theorem}[Inner product of joined blocks]
\lean{FABL.innerProductModTwoBit_joinF₂CubeBlocks}
\leanok
\uses{FABL.f₂DotProduct, FABL.innerProductModTwoBit, FABL.joinF₂CubeBlocks}
\difficulty{1}
For every natural number $n$ and vectors $x,y\in\mathbb F_2^n$, the inner product modulo
two of the vector formed by concatenating $x$ and $y$ is the dot product of $x$ and $y$
in $\mathbb F_2$.
\end{theorem}

\begin{theorem}[Dot product of joined binary blocks]
\lean{FABL.f₂DotProduct_joinF₂CubeBlocks}
\leanok
\uses{FABL.f₂DotProduct, FABL.joinF₂CubeBlocks}
\difficulty{1}
For any positive integer $n$ and vectors $a,b,x,y\in\mathbb F_2^n$, the dot product over
$\mathbb F_2$ of the concatenations of $a,b$ and $x,y$ is the sum of their blockwise dot
products: $\langle(a,b),(x,y)\rangle=\langle a,x\rangle+\langle b,y\rangle$.
\end{theorem}

\begin{theorem}[Sign of a binary inner product]
\lean{FABL.innerProductModTwo_joinF₂CubeBlocks}
\leanok
\uses{FABL.binarySign, FABL.f₂DotProduct, FABL.innerProductModTwo, FABL.innerProductModTwoBit_joinF₂CubeBlocks, FABL.joinF₂CubeBlocks, FABL.realSignEncodedFunction, FABL.signEncodedFunction, FABL.signValue_signEncode_eq_binarySign}
\difficulty{3}
For every natural number $n$ and vectors $x,y\in(\mathbb{F}_2)^n$, the real sign
encoding of their inner product modulo two equals the binary character evaluated at that
inner product: $\operatorname{sign}(\langle x,y\rangle)=\chi(\langle x,y\rangle)$.
\end{theorem}

\begin{theorem}[Walsh characters on concatenated blocks]
\lean{FABL.vectorWalshCharacter_joinF₂CubeBlocks}
\leanok
\uses{FABL.f₂DotProduct_joinF₂CubeBlocks, FABL.joinF₂CubeBlocks, FABL.vectorWalshCharacter, FABL.vectorWalshCharacter_apply}
\difficulty{3}
Let $n$ be a nonnegative integer and let $a,b,x,y\in\mathbb F_2^n$. Concatenate each
pair of vectors to form vectors of length $2n$. The Walsh character indexed by the
concatenation of $a$ and $b$, evaluated at the concatenation of $x$ and $y$, equals the
product of the Walsh character indexed by $a$ evaluated at $x$ and the Walsh character
indexed by $b$ evaluated at $y$.
\end{theorem}

\begin{theorem}[Addition and equality in the Boolean cube]
\lean{FABL.add_eq_zero_iff_eq_f₂Cube}
\leanok
\difficulty{4}
For every natural number $n$ and all $x,y\in\mathbb{F}_2^n$, the sum $x+y$ is zero if
and only if $x=y$.
\end{theorem}

\begin{theorem}[Averaged Walsh character identity]
\lean{FABL.expect_mul_vectorWalshCharacter_add}
\label{FABL.expect_mul_vectorWalshCharacter_add}
\leanok
\uses{FABL.add_eq_zero_iff_eq_f₂Cube, FABL.expect_vectorWalshCharacter, FABL.vectorWalshCharacter}
\difficulty{4}
Let $n$ be a natural number and let $a,b$ be elements of $(\mathbb F_2)^n$. With uniform
expectations over $x,y\in(\mathbb F_2)^n$, the expectation of $\chi_a(x)\chi_{x+b}(y)$
over $x$ and $y$ equals the expectation over $x$ of $\chi_a(x)$ times the indicator that
$x=b$, where $\chi_c$ denotes the Walsh character indexed by $c$.
\end{theorem}

\begin{theorem}[Vector fourier coeff inner product mod two join f₂cube blocks]
\lean{FABL.vectorFourierCoeff_innerProductModTwo_joinF₂CubeBlocks}
\leanok
\uses{FABL.binarySign, FABL.expect_mul_vectorWalshCharacter_add, FABL.f₂CubeBlockEquiv, FABL.f₂DotProduct, FABL.innerProductModTwo, FABL.innerProductModTwo_joinF₂CubeBlocks, FABL.joinF₂CubeBlocks, FABL.vectorFourierCoeff, FABL.vectorFourierCoeff_eq_expect, FABL.vectorWalshCharacter, FABL.vectorWalshCharacter_apply, FABL.vectorWalshCharacter_joinF₂CubeBlocks, FABL.𝔽₂}
\statementsource{Analysis-of-Boolean-Functions-by-Ryan-ODonnell}{pdf-ffbe9c7b3306-p117-b002}
\difficulty{5}
% dataset-pipeline: stub -- prose pending build_informal_statements.py
Statement pending regeneration.
\end{theorem}

\begin{theorem}[Fourier coefficient of the binary inner product]
\lean{FABL.abs_vectorFourierCoeff_innerProductModTwo}
\proofsource{Analysis-of-Boolean-Functions-by-Ryan-ODonnell}{pdf-ffbe9c7b3306-p117-b002}
\label{FABL.abs_vectorFourierCoeff_innerProductModTwo}
\leanok
\uses{FABL.abs_vectorWalshCharacter, FABL.binarySign, FABL.f₂CubeBlockEquiv, FABL.f₂DotProduct, FABL.innerProductModTwo, FABL.joinF₂CubeBlocks, FABL.vectorFourierCoeff, FABL.vectorFourierCoeff_innerProductModTwo_joinF₂CubeBlocks, FABL.vectorWalshCharacter_apply}
\difficulty{4}
For every natural number $n$ and every pair of binary vectors $\gamma$ of length $2n$,
the absolute value of the Fourier coefficient at $\gamma$ of the sign encoding of the
mod-two inner product between the two length-$n$ blocks is $2^{-n}$.
\end{theorem}

\begin{theorem}[Real representation of the inner product function]
\lean{FABL.innerProductModTwoBoolean_toReal}
\label{FABL.innerProductModTwoBoolean_toReal}
\leanok
\uses{FABL.BooleanFunction.toReal, FABL.binaryFunctionOnSignCube, FABL.innerProductModTwo, FABL.innerProductModTwoBoolean, FABL.realSignEncodedFunction, FABL.signEncodedFunction}
\difficulty{3}
For every natural number $n$, the real-valued function obtained by viewing the
sign-valued inner product Boolean function on the sign cube of dimension $2n$ as
real-valued equals the function obtained by transporting the real-valued, sign-encoded
inner product modulo two from the binary cube to the sign cube.
\end{theorem}

\begin{theorem}[Fourier coefficient of the mod-two inner-product function]
\lean{FABL.abs_fourierCoeff_innerProductModTwoBoolean}
\proofsource{Analysis-of-Boolean-Functions-by-Ryan-ODonnell}{pdf-ffbe9c7b3306-p117-b002}
\label{FABL.abs_fourierCoeff_innerProductModTwoBoolean}
\leanok
\uses{FABL.BooleanFunction.toReal, FABL.abs_vectorFourierCoeff_innerProductModTwo, FABL.fourierCoeff, FABL.f₂CubeEquivFinset, FABL.f₂Support, FABL.innerProductModTwoBoolean, FABL.innerProductModTwoBoolean_toReal, FABL.vectorFourierCoeff_eq_fourierCoeff_binaryFunctionOnSignCube}
\difficulty{4}
Let $n$ be a nonnegative integer and let $S$ be any subset of the $2n$ coordinates. The
absolute value of the Fourier coefficient indexed by $S$ of the Boolean function
obtained by encoding the mod-two inner product of the two $n$-bit blocks as a sign is
$2^{-n}$.
\end{theorem}

\begin{theorem}[Pow two le polynomial sparsity inner product mod two]
\lean{FABL.pow_two_le_polynomialSparsity_innerProductModTwo}
\statementsource{Analysis-of-Boolean-Functions-by-Ryan-ODonnell}{
  pdf-ffbe9c7b3306-p117-b001,
  pdf-ffbe9c7b3306-p117-b002,
  pdf-ffbe9c7b3306-p116-b006,
  pdf-ffbe9c7b3306-p116-b007
}
\label{FABL.pow_two_le_polynomialSparsity_innerProductModTwo}
\leanok
\uses{FABL.IsPolynomialThresholdRepresentation, FABL.abs_fourierCoeff_innerProductModTwoBoolean, FABL.fourierSupport, FABL.innerProductModTwoBoolean, FABL.one_le_sum_abs_fourierCoeff_of_polynomialThresholdRepresentation, FABL.polynomialSparsity}
\difficulty{4}
% dataset-pipeline: stub -- prose pending build_informal_statements.py
Statement pending regeneration.
\end{theorem}

\begin{theorem}[Sign encoding of the binary inner product]
\lean{FABL.innerProductModTwoBoolean_binaryCubeSignEquiv_joinF₂CubeBlocks}
\leanok
\uses{FABL.binaryCubeSignEquiv, FABL.f₂DotProduct, FABL.innerProductModTwoBoolean, FABL.joinF₂CubeBlocks, FABL.signEncode}
\difficulty{1}
For any natural number $n$ and vectors $x,y\in\mathbb{F}_2^n$, evaluate the Boolean
function on the sign cube of dimension $2n$ at the sign encoding of the concatenation of
$x$ and $y$. The result is the sign encoding of their inner product modulo two.
\end{theorem}

\begin{theorem}[Nonconstancy of the inner product function]
\lean{FABL.innerProductModTwoBoolean_ne_one_of_pos}
\label{FABL.innerProductModTwoBoolean_ne_one_of_pos}
\leanok
\uses{FABL.binaryCubeSignEquiv, FABL.f₂DotProduct, FABL.innerProductModTwoBoolean, FABL.innerProductModTwoBoolean_binaryCubeSignEquiv_joinF₂CubeBlocks, FABL.joinF₂CubeBlocks}
\difficulty{4}
For every positive natural number $n$, the Boolean function on the sign cube of
dimension $2n$ that sign-encodes the inner product modulo two is not the constant
function with value $1$.
\end{theorem}

\begin{theorem}[Sparsity of inner product thresholds]
\lean{FABL.pow_two_le_polynomialSparsity_innerProductModTwo_of_pos}
\proofsource{Analysis-of-Boolean-Functions-by-Ryan-ODonnell}{
  pdf-ffbe9c7b3306-p116-b006,
  pdf-ffbe9c7b3306-p116-b007,
  pdf-ffbe9c7b3306-p117-b001,
  pdf-ffbe9c7b3306-p117-b002
}
\label{FABL.pow_two_le_polynomialSparsity_innerProductModTwo_of_pos}
\leanok
\uses{FABL.IsPolynomialThresholdRepresentation, FABL.innerProductModTwoBoolean, FABL.innerProductModTwoBoolean_ne_one_of_pos, FABL.polynomialSparsity, FABL.pow_two_le_polynomialSparsity_innerProductModTwo, FABL.thresholdSign}
\difficulty{4}
Let $n>0$, and consider the Boolean function on pairs of $n$-bit strings that outputs
their dot product over $\mathbb F_2$, encoded as a sign. If a real-valued function on
the corresponding sign cube represents this function by threshold signs, then it has at
least $2^n$ nonzero Fourier coefficients.
\end{theorem}
